\section{Items, records, and inferential targets}\label{sec:estimand}
There are $k$ distinct items, with reference labels $y_j\in\{0,1\}$ and
predictions $q_j\in\{0,1\}$. A known $n\times k$ matrix $C_r$ maps them into
the $r$th record-level measure:
\[
T_r=C_ry,\qquad S_r=C_rq,\qquad d_r=S_r-T_r=C_re,\qquad e=q-y.
\]
Rows can represent people, places, organizations, or periods; columns
represent the labeled items. Entries can be counts, durations, expenditures,
or observed shares. An item may contribute to several rows and several
matrices. Different aggregates therefore share the same vector $e$.

For fixed record weights $a$, a linear target and its prediction-based
estimate are
\begin{equation}\label{eq:estimand}
\theta=a^{\mathsf T}C_ry=w^{\mathsf T}y,\quad
\widehat\theta=w^{\mathsf T}q,\quad w=C_r^{\mathsf T}a.
\end{equation}
Totals have nonnegative weights; comparisons can have signed weights.
Proportions with observed denominators fit this expression after the
denominators are absorbed into $C_r$. A denominator depending on unknown
labels, a regression using $C_rq$ as a covariate, or another nonlinear
transformation instead has error $F(C_1q,\ldots,C_Rq,Z)-F(C_1y,\ldots,C_Ry,Z)$,
where $Z$ contains other observed inputs. For linear targets the exact error is
\begin{equation}\label{eq:contrast}
\Delta=\widehat\theta-\theta=w^{\mathsf T}e.
\end{equation}

For illustration, consider three records and two items with
\[
C=\begin{pmatrix}4&0\\1&2\\0&1\end{pmatrix},\qquad
y=(0,1)^{\mathsf T},\qquad q=(1,1)^{\mathsf T}.
\]
One false positive changes the record measures from $T=(0,2,1)^{\mathsf T}$
to $S=(4,3,1)^{\mathsf T}$. With $Z=(0,1,2)^{\mathsf T}$ and an intercept,
regressing $T$ on $Z$ gives slope $1/2$; substituting $S$ gives $-3/2$.
Regressing $Z$ on $T$ also gives $1/2$, while regressing $Z$ on $S$ gives
$-9/14$. The same item mistake has different consequences in the two roles.

The primary target below is a descriptive coefficient for this observed
panel with complete item labels. The matrices, observed regressors, and
observed outcomes are held fixed when calculating its label error.
Uncertainty requires an additional repeated experiment. A prediction-error
model varies $q$ through a specified fitting or prediction process while
holding $y$ fixed. An item-validation design instead varies which fixed
labels are verified while holding $q$ fixed. Section~\ref{sec:inference}
uses the latter to construct feasible corrections. Neither experiment by
itself establishes inference about a population or causal coefficient.

If records are sampled as well, their sampling distribution must be stated.
An SE conditional on frozen predictions concerns that record-sampling
problem. It need not quantify uncertainty about the complete-label
coefficient. The corresponding variance components can be dependent; adding
a label-error variance to a reported regression variance is justified only
after deriving their joint distribution.
